\section{Common windows for the prime pairs and the Selmer estimate}
\label{sec:rh-common-windows}

Fix a primitive character $\psi$ of conductor $m$, $0<\epsilon<1/52$,
and a scale factor $c\ge1$. Put $z=U^{\epsilon/5}$,
$z'=z(\log U)^2$, $P=\exp z$, and $D=\exp z'$. For every fixed $A\ge1$
all sufficiently large $U$ admit primes $q,d$ with $P^{1/2}<q\le P$,
$D^{1/2}<d\le D$, and $(m,qd)=1$. They satisfy $q\equiv3\bmod 8$,
$d\equiv7\bmod16$, $q\mid d+1$, and $p\mid d+1$ for every odd prime
$p\le z$. For every character $\xi$ modulo $m$, every nonzero
$(i,j)\in\{0,1\}^2$, and every primitive character $\chi$ inducing
$\xi(\frac{\cdot}{q})^i(\frac{\cdot}{d})^j$ on units prime to $2mqd$,
this same choice gives
\[
 |G_\chi(y)|\le y^{1/2+\epsilon/2},\qquad
 |L'(1,\chi)/L(1,\chi)|\le4\log y\qquad (U\le y\le U^A).
\]
Here $G_\chi$ is the complex centered prime sum. For any primitive real
$\rho$ with conductor dividing $8m$ in place of $\xi$, the corresponding
inducing character has no zero for $\Re s>1-1/400$ and
$|\Im s|\le8mPD^j$. These are the conclusions of \texttt{PairGood}.

\subsection{Simultaneous choice of the primes}

The first candidates are the primes in $(\sqrt P,P]$ congruent to
$3\bmod8$. For each candidate $q$, the second candidates are the primes
in $(\sqrt D,D]$ congruent to $7\bmod16$ and to $-1$ modulo $q$ and
all odd primes at most $z$. These congruences form a single reduced
residue class modulo $16q\prod_{2<p\le z}p$ for all sufficiently
large $U$, with the product taken over primes. For every fixed
$\eta>0$ the estimates \texttt{stage1} and \texttt{stage2} exclude
at most an $\eta$ fraction in each selection. The second estimate is
uniform in the first candidate. Both candidate sets are nonempty.
Taking $\eta=1/2$ leaves a $q$ meeting every first-stage requirement,
then a $d$ meeting every requirement for that $q$. In the first stage
the family comprises the primitive representatives of all characters
modulo $8m$. In the second these characters are multiplied by
$(\frac{\cdot}{q})^i$, $i=0,1$, before taking primitive representatives;
the number of characters to test is at most twice the number modulo
$8m$. For large $U$, $d>q>8m$.

Here are the exceptional-set bounds in \texttt{StagePow}. Write $V$
for the upper endpoint of the candidate progression, $\mathcal P$ for
the candidate set, and $F$ for the maximum allowed base conductor.
Then $1\le F\le\sqrt V$, each candidate exceeds $F$, and the
prime-progression lower bound in the admissible ranges gives
$|\mathcal P|\ge V^{9/10}$. For $Q\ge3$ the zero-density input
\texttt{CountZAll.zero\_density\_box} bounds the number of primitive
characters of conductor at most $Q$ with a zero in
$\Re s>1-1/400,\ |\Im s|\le Q$ by $C Q^{1/4}$.
For a failing prime $p$, choose one bad primitive twist. Its conductor
equals $fp$ with $1\le f\le F$. Distinct primes yield distinct
characters: equality gives $fp=f'p'$, hence $p\mid f'p'$ and
$p=p'$ since $0<f'<p$. Taking $Q=FV$ bounds the number discarded by
$CV^{3/8}$, at most $(\eta/2)|\mathcal P|$ once $V$ is large.
The zero-density theorem imports mollifier zero detection, a block
large-sieve estimate and a coefficient-square bound.

The second exceptional set concerns a bad value at some real point,
which can depend on the prime. With a fixed truncation $T$ put
\[
  S_p(y)=\sum_{n\le T}\Lambda(n)\min(1,y/n)\kappa(n)
                             (\tfrac np),\qquad |\kappa(n)|\le1 .
\]
The mean-square input \texttt{meanSqK} gives
$\sum_{p\in\mathcal P}|S_p(y)|^2\le C y\log y\,|\mathcal P|$
uniformly in the permitted $y$ and in $\kappa$. Let $D_p(y)$ denote
the piecewise derivative in $y$. Then $yD_p(y)$ equals $S_p(y)$ with
$\kappa(n)$ replaced by $\kappa(n)\mathbf 1_{n>y}$ away from its
finitely many corners. Thus $\sum_p|D_p(y)|^2\le C(\log y)|\mathcal
P|/y$. For $L\le W\le2L$, integration and Cauchy--Schwarz give
\[
  \sup_{L\le y\le W}|S_p(y)|^2
      \le 2|S_p(L)|^2+2L\int_L^W|D_p(v)|^2\,dv.
\]
Sum over $p$ and apply the two mean-square bounds. The result is
$O(CL\log(2L)|\mathcal P|)$. Markov's inequality bounds the fraction
with $|S_p(y)|>y^{1/2+\tau}/2$ at some point in the block by
$O(C\log(2L)L^{-2\tau})$. For any fixed $\tau>0$ the explicit
relaxation used in \texttt{BadShare} at large $L$ is
$(32C/\tau)L^{-\tau}$. Summing over blocks starting at $2^hU$ bounds
the total fraction by $(32C/\tau)U^{-\tau}/(1-2^{-\tau})$, uniformly
in the upper endpoint. For at most $R$ multipliers, take the last
fraction at most $\eta/(2(R+1))$ and use the union bound. Including
the zero-density exceptions gives at most $\eta|\mathcal P|$ in total.
\texttt{BoxOfS} supplies the explicit-formula comparison to the twists.
Take $\tau=\epsilon/2$. The constant four for the logarithmic derivative
comes from \texttt{kappa4} at heights $H=8mPD^j$ satisfying
$U^{\epsilon/6}\le\log H\le3U^{\epsilon/5}(\log U)^2$.
The prime-progression, mean-square, zero-density and explicit-formula
inputs cited here have separate proofs in the source.

\subsection{Transferring both estimates to one interval}

Take $a(u)$ with $u^{2+8\epsilon}\le2^{a(u)}\le2u^{2+8\epsilon}$ for
$u\ge1$. Apply the strict window theorem to obtain $J\ge1,X_0$, and
the good-pair theorem with $A=J$ to obtain $U_0$. Set
$k=c^{1/(2+8\epsilon)}$. Given $y_0$, choose
$M\ge\max(1,y_0,U_0,k^3,k^{400})$ and $X\ge\max(X_0,M^J)$.
Set $U=X^{1/J}$. Choose one good pair at $U$ through $U^J=X$.
Apply the strict window theorem to this $X,q,d$. It supplies $W\ge U$
with $W^2\le X$ and the strict estimates throughout $[W,W^2]$.

Put $Z=kW$. Since $k^3\le W$, we have $Z^{3/2}\le W^2$. For every
$x\in[Z,Z^{3/2}]$ the point $u=x/k$ satisfies $W\le u\le x\le W^2$.
Set $e=a(u)$. Then $x^{1/J}\le U\le x$ and the same $q,d$ work for
every argument from $U$ to $x$. Furthermore
\[
 x^{2+8\epsilon}\le c2^e\le2x^{2+8\epsilon},
 \qquad x^{400}=k^{400}u^{400}\le u^{401}.
\]
For every field $E$ realizing the specified fixing subgroup for
$(\psi,e,q,d)$, apply the strict estimate at $u$, shorten the range
of twists, and increase the right side using $u^\epsilon\le x^\epsilon$.
This gives $\mathrm{VVle}(E,\lfloor x^{400}\rfloor,
\sqrt{\exp(x^\epsilon)})$ as in \texttt{WinPairs.winPairs}.
Here $Z\ge y_0$ and $q,d$ stay fixed on the whole indicated interval.

\subsection{The family used in the analytic ending}

The level $\Lambda$ consists of the primitive representatives of the
powers of $\psi$, including $\psi$ and the primitive trivial character.
The weight is one for the trivial member and the distinguished member
if present, and $1/(3\Lambda.o+1)$ for every other member. Thus the
weighted centered sum $G_\Lambda$ is fixed once $\psi$ is fixed. The
host requirement $G_\psi$ asks for the character comparisons
$\mathrm{EAtH}(\Lambda,H,T,x,\epsilon,4)$ together with the strict
bound $\mathrm{VVle}(E,t,\sqrt b)$ on that same host. In
\texttt{supplyOfAdmOf} the admissible-rank construction yields the
host. The prime rows and twist estimates give the first conjunct.
The fixing subgroup identifies its field with the field of the strict
estimate, and $t\le x^{400}$ gives the second conjunct. The resulting
host supply \texttt{supplyChi} has window exponent $\lambda=3/2$.

In \texttt{topDesWin} the closed estimates for the level, divisibility,
auxiliary character, irreducibility, small-level rank, selected rank,
and the three comparisons $G,S,\mathrm{GLOBAL}$ are conjoined. Their
conjunction gives a positive $\epsilon_1$ and, for each sufficiently
small fixed $\epsilon$, one lower-bound tuple $K$ for all nine.
Design existence gives a positive $\epsilon_3$ and quantifies over
every such $K$. Take $\epsilon<\min(\epsilon_1,\epsilon_3)$ and
select $K$ before selecting the design family. On the resulting
family its two size bounds, level and character comparisons and
strict Selmer estimate imply the divisibility relation. Its rank
conclusions then give the required multiplicity relation and the
global comparison supplies the chain. In \texttt{shapeWin} the
specific exponents passed into this argument are $e_D=2,e_L=0$.
The nine comparison inputs are imported here.

Consequently for some $\epsilon_0>0$ and every $0<\epsilon<\epsilon_0$
there is a family with a scale set $S$ containing arbitrarily late
intervals $[Z,Z^{3/2}]$. Beyond its threshold every host index
satisfies the same character comparison with constant four. On its
sufficiently large scales in $S$ the compatible chain gives, by
\texttt{tiedCore\_of\_chainTied}, a host index witnessing the core
numerical comparison. The character comparison applies to that index
because it applies to all indices. Beyond the maximum of the three
thresholds for the chain, the character comparison and
\texttt{core\_gm}, the latter estimate yields
\[
        G_\Lambda(x)\le x^{1/2+(e_D+e_L+7)\epsilon}\qquad(x\in S).
\]
For $\Re s>1/2$, put $b=\tfrac12\min(\Re s-1/2,1/2)$ and
$C=e_D+e_L+7$. With $\epsilon=\min(\epsilon_0/2,b/(|C|+1))$ one
gets $C\epsilon\le b$ and $1/2<\theta=1/2+b<\min(1,\Re s)$.
Given any lower endpoint, choose an entire interval above it, above
one, and beyond all three thresholds. Throughout it $G_\Lambda(x)
\le x^\theta$. The analytic window criterion now gives
$L(s,\psi)\ne0$. This is the deduction in
\texttt{zeroFreeAll\_of\_topG} from \texttt{level\_of\_tiedWin}.
